% !TEX root = ../main.tex
%Policy makers choice over grading regimes decides how various groups of students gets access to universities with varying selectivity levels, with broader implications for the role of universities on social mobility.
This paper tests how a policy-maker's choice over grading regimes changes the allocation of students across universities by using the cancellation of standardized tests in the UK in 2020-2021. The UK replaced centralized exams (e.g. A-levels) with teacher-assessed grades after universities had already made conditional offers based on expected exam results. This allowed each school to set their own grading leniency levels that determined the application outcomes of their students. As a result, the grading regime switch pushed students with low predicted A-levels attainments into universities that they would not have been accepted otherwise. The setting allows me to the study whether the compositional shifts were proportional to the applicants schools' choice over grading leniency. 

I show that schools exercised their grading discretion to very different degrees. Schools serving economically affluent students adopted looser grading standards, increasing their students' representation at elite university programs. I identify the grading discretion exerted at each school by measuring the deviations in the average student's grade on years that the A-levels were cancelled, and comparing them to the school's historical levels, which were mechanically fixed to follow a common distribution by the interventions of the central graders. I then estimate how the influx of students with inflated grades changed the composition of university programs, using the average grade inflation of each program's applicant pool as an instrument.  

%Results: I provide 3 sets of new results.
I show private institutions (e.g. \emph{independent schools}) were more likely to inflate their student's grade, and rapidly accerating the degree of inflation in the second year.  Private schools increased their students chances of obtaining a top grade (e.g. A/A*) by 8 percentage more than other publicly funded schools. Additionally, indepedent schools adjusted their grading standards to a much looser level in the second year, as the attainment gap between independent schools and publicly run school increased by roughly 50 percent in the second year. The impact of the grade increments on placement success was much larger for students applying to selective universities, and the share of students in the incoming cohorts increased by [] percentage point relative to average levels in normal years, while reducing the share of students from publicly funded schools by [] percentage points. Within independent schools, schools with higher value added showed a large increase in their students final grades.

%Results 2: Schools inflated through non-STEM subjects. 
I show how grading discretion was strongly observed in non-STEM subjects and within-course differences by gender, race, and SES.  I find looser grading policies were uniformly observed in non-STEM subjects, as students incremental chances of obtaining a top grade was nearly twice as high as academically similar students in STEM subjects. Within a subject course, schools assigned higher grades to female, white, and high SES students. Additionally, the gap was pronouned for independent schools as the gap between gender, race, and high SES was []\% higher than their publicly funded school counterpart. 

%Results 3: Universties responded by strictening grades, but it excluded students from less advantaged backgrounds.
Lastly, I show the grade inflation were enough to shift the ordering of student within the applicant pool, and students from schools that did not inflate their students got reject more likely when universities responded to the grade regime change by strictening their grade requirements. A one unit increase in the grade requirement of an elite university caused a []\% and [] \% reduction in a students from a publicly funded school and independent school to get accepted into the program, which was []\% and []\% during non-COVID years. 

% Motivating the Contribution
The inflow of students into universities caused by the transfer of grading discretion, triggered responses from universities in tightening their entry requirements under constraints\footnote{UK universities set uniform entry conditions to every applicant} and schools responded in the second year by revising the grading standards. The first year grade inflation increased the number of accepted students, especially at selective institutions. Universities with selectivity levels in the top 10 percentile groups expanded their incoming cohort size by nearly 60 \% in 2020. The expansion exhausted each courses teaching resources, and many programs tightened their admission standards in the second year. Despite the efforts, the schools further reduced their grading standards in the second years, and leading to an even larger increase in the incoming student size in the second year. 

% Contribution 1: Evidence of how admission policies affect student composition through dynamic adjustments.
By tightening entry requirements to limit enrollment, the university replaced high-achieving, socio-economically disadvantaged students with lower-achieving students from wealthy families, displacing the disadvantaged students to lower-ranked institutions. Previous papers highlight
how internally choosen admission policies (\cite{Chetty2023}) or externally imposed rules (\cite{BLEEMER2023104839}) determines the groups of students entering elite universities, relagating unselected students into lower ranked university programs. However, schools responds to such policy changes by altering the applicant pool by modifying the qualifications presented to universities. As a result, the combination of the admission policy change and schools responses in assining different qualifications to students may result in sub-optimal outcomes for both universities and the students. 

% Constribution 2: Evidence of grade inflation
Beyond the primary enrollment effects, my setting reveals how schools translate grading discretion into shifts in the relative ordering of students along the qualification distribution. Schools exercised this discretion to very different degrees, and because students were non-randomly sorted across schools and subjects before the policy was announced, the resulting grade jumps differed sharply across demographic groups. Two channels drove these differences. First, grade revisions were substantially larger in non-STEM subjects (\cite{Ahn2024}), so students whose subject portfolios tilted toward the humanities experienced larger upward shifts. Second, within a given subject-course, schools assigned differentially higher grades to students from advantaged backgrounds — a pattern consistent with teacher bias in subjective assessment (\cite{hanna2012}; \cite{Alesina2024}) and with differences in the unobserved attributes that teachers reward (\cite{cornwell2013}). Importantly, schools varied not only in the overall level of inflation they applied but also in how much their grading distinguished across demographic groups: independent schools and schools with higher value-added scores exhibited larger within-school gaps in grade revisions by gender, race, and socioeconomic status. Finally, these disparities compounded over time. As schools recalibrated and universities tightened entry requirements in the second year, the pre-existing gaps across student groups widened further, amplifying the initial sorting advantage of students in schools that inflated aggressively.

The paper is organized as follows: Section 2 provides the institutional background of the higher education system in the UK as well as the details of the policy set in place during the pandemic. Section 3 describes the data source I use in the analysis. Section 4 presents the empirical model and discuss the validity of the estimates along with the institutional details. Section 5 presents the results. Section 6 discuss the results with the literature. Section 7 concludes with policy discussions. 